\clearpage
\fancyhead[L]{\small HALLASCHKA. ELEMENTA ECLIPSIUM.}
\SourceLeaf{624}
\begin{center}\small Göttingische gelehrte Anzeigen. No. 17. Pages 167..168. 29 January 1820.\\[3pt]\rule{.72\linewidth}{.3pt}\end{center}

\emph{Elementa eclipsium quas patitur tellus, luna eam inter et solem versante, ab a. {\upshape 1816} usque ad a. {\upshape 1860}, ex tabulis astronomicis recentissime conditis et calculo parallactico deducta, typo ecliptico et tabulis projectionis geographicis collustraia, a} \textsc{Cassiano Hallaschka} Prague 1816. Typis Theophili Haase. 107 pages in quarto, with 20 copperplates.

\begingroup\brokenpenalty=10000\relax
By far the largest part of this work is taken up by a guide to the parallactic calculation and orthographic plotting of solar eclipses. The guide uses the solar eclipses of 1816 and 1818 as practical examples and gives the calculations for the former in full detail. For complete beginners, these instructions, based solely on the usual simple methods, may be of some use. With these readers in mind, however, the author should have taken somewhat greater care to present the preliminary directions clearly and accurately. On the last 18 pages the author also gives the principal circumstances of the eclipses visible in Brünn from 1818 to 1860, while the accompanying copperplates give a general representation of those eclipses for the whole Earth. Among them are three partial, five annular and four total eclipses. For the horizon of Brünn, the eclipse of 8 July 1842 is the greatest, and the magnitude of the eclipse there is 11 digits $13'$. The drawing itself does not seem to have been made with the greatest care throughout. In the eclipse just mentioned, we find that the points at which it begins or ends at sunrise or sunset are drawn as a single continuous line, though they necessarily had to form two separate lines joined by the line of contact between the Sun’s southern limb and the Moon’s northern limb. The calculations were carried out using \textsc{Triesnecker}'s lunar tables in the Vienna Ephemerides for 1803 and manuscript solar tables by the same astronomer. The prime meridian is set, instead of the usual $20^\circ$, at $19^\circ\,53'\,45''$ west of Paris.
\par\clearpage\endgroup
